\section*{Declarations}

\bmhead{Ethics approval and consent to participate}
Not applicable.

\bmhead{Consent for publication}
Not applicable.

\bmhead{Availability of data and materials}
\textbf{Software availability.}
\begin{description}
\item[Project name:] QSARena
\item[Project home page:] \url{https://github.com/ScottCoffin/QSARena}
\item[Operating system(s):] Platform independent (Windows, macOS, Linux, Google Colab)
\item[Programming language:] Python 3.11
\item[Installation:] \texttt{pip install qsarena} (console scripts \texttt{qsarena-benchmark}, \texttt{qsarena-applicability-domain} and \texttt{qsarena-examples}), or clone the repository and run the scripts directly; the code-free notebook requires no installation at all and runs in Google Colab. A step-by-step installation and usage tutorial, in which every command is executed by an automated test on bundled example data, is provided as Additional file~2
\item[Other requirements:] RDKit, scikit-learn, pandas, NumPy, SciPy, joblib, PyYAML; heavy backends are optional extras (\texttt{qsarena[boosting]}, \texttt{[deep]}, \texttt{[graph]}, \texttt{[foundation]}, \texttt{[notebook]}). Exact versions for reproducing the benchmark are pinned in \texttt{requirements-cpu.txt} / \texttt{requirements-cuda.txt} and the conda environment files, which should be used in preference to the pip extras for exact reproduction
\item[License:] MIT License (OSI-approved); see the \texttt{LICENSE} file in the repository
\item[Any restrictions to use by non-academics:] None beyond the repository license
\end{description}

The QSARena source code, benchmark runner, analysis notebook and the complete benchmark artifacts analysed in this study are available in the QSARena repository, \url{https://github.com/ScottCoffin/QSARena}. The repository contains the workflow core (\texttt{qsar\_workflow\_core.py}), the notebook builder, the command-line benchmark runner (\texttt{run\_qsarena\_benchmarks.py}, installed as the \texttt{qsarena-benchmark} console script), the dataset registry, the curated leaderboard reference tables, the per-dataset benchmark artifacts (metrics, selected features, selector coefficients, split-signature hashes, runtime records, fusion candidate tables and ensemble weights) for the run from which every reported result is computed (\texttt{benchmark\_results/qsarena\_benchmark\_oof\_ensemble/}) and for the runs it was built from (\texttt{benchmark\_results/autoqsar\_benchmark\_20260623\_153839/} and \texttt{benchmark\_results/qsarena\_benchmark\_chemprop\_fixed/}) together with the consumer-GPU (RTX 4060 laptop) comparison run of Section~\ref{sec:hardware}, whose directory is named in Additional file~1, Note~S1, the post-hoc feature-expansion analysis of Additional file~1, Note~S3 (\texttt{benchmark\_results/qsarena\_feature\_expansion/}), the analysis notebook (\texttt{portable\_colab\_qsar\_bundle/benchmark\_results\_summary.ipynb}), and the scripts that regenerate every figure, table and quoted number in this paper (\texttt{render\_manuscript\_assets.py}, \texttt{verify\_manuscript\_numbers.py}). Pinned conda and pip environment specifications, an Apptainer container definition and Slurm/OpenStack submission scripts are included for independent re-execution.

The datasets analysed during the current study are available in the following public repositories: the Therapeutics Data Commons ADMET Benchmark Group and single-prediction tasks, via PyTDC (\url{https://tdcommons.ai}) \cite{huang2021tdc}; the MoleculeNet ESOL \cite{delaney2004esol}, FreeSolv \cite{mobley2014freesolv} and Lipophilicity (ChEMBL deposition CHEMBL3301361, \url{https://doi.org/10.6019/CHEMBL3301361}) \cite{astrazeneca2015chembl} datasets, as distributed by MoleculeNet \cite{wu2018moleculenet}; the Polaris ADME benchmarks (\url{https://polarishub.io}) \cite{ash2025polaris,wognum2025polarissoftware}, derived from the Biogen ADME dataset \cite{fang2023adme}; the PODUAM point-of-departure datasets (\url{https://github.com/kejbo/PODUAM}) \cite{vonborries2026uncertainty}; and the ChemML example datasets \cite{haghighatlari2020chemml}. Each dataset's source, version and split are recorded in the dataset registry and in Additional file~1, Table~S6.

No registration or login is needed to download, install or run QSARena, so reviewers can test it anonymously. The optional Google Colab path requires a Google account; the same notebook also runs in any local Jupyter installation, and Additional file~2 covers local installation step by step.

A versioned release archived on Zenodo, including the per-molecule prediction files excluded from the repository for size, will be deposited before publication and its DOI cited here.
% TODO(author): mint Zenodo DOI and replace this sentence. The deposit is staged: .zenodo.json and CITATION.cff hold the
% metadata, and dist/zenodo/ holds the two per-molecule prediction archives with sha256 manifests (re-verified against
% the repository files on 2026-10-06; ZENODO.md gives the steps).

\bmhead{Competing interests}
The author declares no competing interests.

\bmhead{Funding}
This work was supported by the California Office of Environmental Health Hazard Assessment (OEHHA). Computing resources for the benchmark were provided in kind through NSF ACCESS allocation CIS261142 on Jetstream2 (see Acknowledgements). The funders had no role in study design, data collection and analysis, decision to publish, or preparation of the manuscript.

\bmhead{Authors' contributions}
SC conceived the study, developed the software, designed and executed the benchmark, analysed the results and wrote the manuscript. The author read and approved the final manuscript.

\bmhead{Acknowledgements}
This work used Jetstream2 GPU at Indiana University through allocation CIS261142 from the Advanced Cyberinfrastructure Coordination Ecosystem: Services \& Support (ACCESS) program, which is supported by U.S. National Science Foundation grants \#2138259, \#2138286, \#2138307, \#2137603, and \#2138296 \cite{boerner2023access}. Jetstream2 \cite{hancock2021jetstream2} is supported by the National Science Foundation under Grant 2005506. Any opinions, findings, and conclusions or recommendations expressed in this material are those of the author and do not necessarily reflect the views of the National Science Foundation.

The views expressed are those of the author and do not necessarily represent those of the California Environmental Protection Agency or the Office of Environmental Health Hazard Assessment.
% TODO(author): confirm the required OEHHA/CalEPA disclaimer wording.
